\documentclass[gmd, manuscript]{copernicus}

\begin{document}

\begin{abstract}
%% DRAFT BULLETS (~250 words when written out):
\begin{itemize}
\item Problem: the layout that makes data FAIR is not the layout that trains fast.
\item Two representations: interoperable Zarr store vs.\ training-optimised merged store.
\item Result: $\sim$7$\times$ per-sample, $\sim$4.6$\times$ end-to-end, $\sim$40$\times$ projected.
\item Contribution: an automated bridge, plus the released 1850--2100 land-use cube.
\end{itemize}
\end{abstract}
\copyrightstatement{TEXT} %% This section is optional and can be used for copyright transfers.
%% ==================================================================
%% 1  INTRODUCTION
%% ==================================================================
\introduction  %% \introduction[modified heading if necessary]
\begin{itemize}
\item Context: dataloading is now the bottleneck for ML on Earth-system data at HPC scale.
\item Problem: FAIR stores and fast stores pull in opposite directions.
\item Gap: the trade-off is unquantified and no bridge exists.
\item Setting: AI4Land pipeline on MareNostrum 5.
\end{itemize}

\subsection{The interoperability--efficiency tension for ML-ready geospatial data}
\begin{itemize}
\item The two poles: reusability vs.\ throughput.
\item Why they conflict: many small objects stress the metadata server.
\item Claim: general to any chunked format, not specific to this dataset.
\end{itemize}

\subsection{Related work}
\begin{itemize}
\item ARCO/Zarr storage for Earth-system data.
\item I/O optimisation on HPC; GPFS/Lustre metadata studies.
\item FAIR principles and ML dataset metadata (Croissant).
\item Dataloading formats at scale (WebDataset, TFRecord).
\item Our position: the intersection --- trade-off plus tooling.
\end{itemize}

\subsection{Contributions and paper outline}
\begin{itemize}
\item C1: quantified trade-off, with a reusable diagnostic methodology.
\item C2: an automated bridge between representations.
\item C3: the released, DOI'd 1850--2100 cube.
\end{itemize}

%% ==================================================================
%% 2  THE DATASET SUBSTRATE
%% ==================================================================
\section{The dataset substrate: a harmonised 1850--2100 land-use cube}
\begin{itemize}
\item What it is: 1850--2100, seven sources, common 1\,km grid.
\item Harmonisation: cite Paper 1, do not re-derive.
\item Role here: the substrate the study operates on and the artefact released.
\end{itemize}

\subsection{Source datasets and provenance}
\begin{itemize}
\item Target labels: HILDA+ v1.0 \citep{winkler2020hilda+}, distributed via PANGAEA.
\item Dynamic inputs: LUH2 v2h/v2f \citep{hurtt2020harmonization} obtained through input4MIPs/ESGF;
      K{\"o}ppen--Geiger climate classes \citep{beck2018present, beck2023high} from GLoH2O;
      and gridded population from the ISIMIP HistSoc protocol \citep{frieler2017assessing}.
\item Static inputs: GEBCO 2024 elevation \citep{GEBCO2024}, together with the PNNL geophysical
      (slope, aspect) and soil (texture, organic carbon) parameter sets \citep{li2024global}.
\item All sources are cited at their original repositories (PANGAEA, UMD/ESGF, GEBCO, GLoH2O,
      the ISIMIP repository, and the PNNL DataHub); CERISE is a delivery channel and appears
      only in the acknowledgements.
\item Version note: the pipeline uses HILDA+ v1.0; v2.0 (2025) was released after the pipeline
      was finalised, and versions are kept consistent with \citet{mozaffari2026scalable}.
\end{itemize}

\subsection{Spatiotemporal coverage and grid}
\begin{itemize}
\item Temporal: back-cast, observed core, SSP forward extension.
\item Spatial: 1\,km grid; extent, CRS, regridding (table).
\item Variables: type, units, temporal resolution (table).
\end{itemize}

\subsection{Licensing and release}
\begin{itemize}
\item All seven sources redistributable with attribution; no share-alike or NC terms.
\item Derived cube released CC BY 4.0 with a DOI (repository TBC).
\item Per-source terms in Table~\ref{tab:licences}.
\end{itemize}

%t
\begin{table*}[t]
\caption{Upstream sources of the released cube, their role, licence, and
redistribution status after modification.}
\label{tab:licences}
\begin{tabular}{llll}
\tophline
Dataset & Role & Licence & Redistribute \\
\middlehline
HILDA+ v1.0 \citep{winkler2020hilda+}        & Target labels             & CC BY 4.0                  & Yes \\
LUH2 v2h/v2f \citep{hurtt2020harmonization}  & Dynamic input             & CC BY 4.0 (CMIP6 terms)    & Yes \\
K{\"o}ppen--Geiger \citep{beck2018present}   & Dynamic input             & CC BY 4.0                  & Yes \\
ISIMIP HistSoc population \citep{frieler2017assessing} & Dynamic input   & CC BY 4.0 (ISIMIP2/2b)     & Yes \\
GEBCO 2024 \citep{GEBCO2024}                 & Static input (elevation)  & Public domain, attribution & Yes \\
PNNL geophysical \citep{li2024global}        & Static input (slope, aspect) & CC0                     & Yes \\
PNNL soil \citep{li2024global}               & Static input (texture, organic) & CC0                  & Yes \\
\bottomhline
\end{tabular}
\belowtable{CMIP6 terms for LUH2 were relaxed to CC BY 4.0 in October 2022.
CC0 is the PNNL DataHub default.}
\end{table*}


%% ==================================================================
%% 3  THE TWO REPRESENTATIONS FAIR <-> FAST
%% ==================================================================
\section{Two representations of one dataset}
\begin{itemize}
\item One cube, two on-disk forms, two different consumers.
\item Different objects, not different views: layout, dtype, metadata all differ.
\item Figure: side-by-side schematic of the two layouts.
\end{itemize}

\subsection{The interoperable (FAIR) store}
\begin{itemize}
\item One array per variable, rich metadata, self-describing.
\item Readable with standard tooling; consolidated metadata for discovery.
\item The published, DOI'd artefact; optionally carries Croissant.
\end{itemize}
\subsection{The training-optimised merged store}
\begin{itemize}
\item Variables stacked into one array, chunked to the sampling pattern.
\item Categoricals stored as float, cast back at load time.
\item Needs a second on-disk copy; fast because it discards self-description.
\item Table: chunk shape, dtype policy, compression, object count.
\end{itemize}
\subsection{What each representation gives up}
\begin{itemize}
\item FAIR store gives up throughput.
\item Merged store gives up semantics, general usability, and storage.
\item Sets up the quantification in Sects.~4 and~6.
\end{itemize}
%% ==================================================================
%% 4  THE EFFICIENCY STUDY  - Ferran's work 
%% ==================================================================
\section{Diagnosing the I/O bottleneck}
\begin{itemize}
\item Training is input-bound, not compute-bound --- but which layer?
\item Approach: isolate dataloader from training, then filesystem from dataloader.
\item Methodology here; numbers in Sect.~6.
\end{itemize}
\subsection{Diagnostic methodology}
\begin{itemize}
\item Two instruments, both reusable beyond this dataset.
\item Measure under realistic shuffled access, not sequential streaming.
\end{itemize}
\subsubsection{Standalone dataloader profiler}
\begin{itemize}
\item Replays the training access pattern without the model.
\item Decomposes per-sample latency; sweeps workers, batch size, store.
\item Separates cold- and warm-cache behaviour.
\end{itemize}
\subsubsection{Synthetic filesystem probe}
\begin{itemize}
\item Controlled open/read patterns on synthetic file populations.
\item Isolates metadata latency from data bandwidth.
\item Establishes: file count $\rightarrow$ metadata ops $\rightarrow$ latency.
\end{itemize}
\subsection{Experimental setup: MareNostrum 5}
\begin{itemize}
\item Hardware: node configuration, partition, memory.
\item Storage: GPFS; metadata service and client caching.
\item Stores tested: 66\,GB and 213\,GB, plus the multi-Zarr baseline.
\item Software stack, pinned environment, benchmark harness.
\end{itemize}
\subsection{Findings}
\begin{itemize}
\item F1: file count and metadata latency dominate, not bandwidth.
\item F2: cache-miss rate scales with store size.
\item F3: store layout is the largest single lever on training speed.
\item Principle: transferable to any chunked-store ML workload on HPC.
\end{itemize}
%% ==================================================================
%% 5  THE BRIDGE 
%% ==================================================================
\section{An automated bridge between representations}
%% THE INTELLECTUAL NOVELTY. Claim depends on the plan §6.2 decision:
%% one-directional (FAIR -> training) vs. bidirectional (round-trip).
\begin{itemize}
\item Beyond ``we made training faster'': users should not have to choose.
\item Publish FAIR; derive the training store locally and reproducibly.
\item Scope: one-directional, with round-trip as demonstrated feature or future work.
\end{itemize}
\subsection{Conversion pipeline design}
\begin{itemize}
\item Read $\rightarrow$ stack $\rightarrow$ re-chunk $\rightarrow$ write; parallelised for HPC.
\item Configuration: chunk shape, dtype policy, variable selection, compression.
\item Manifest (source DOI, version, checksums) makes every derived store traceable.
\item Conversion cost quantified in Sect.~6.4.
\end{itemize}
\subsection{Metadata handling and reconstruction}
\begin{itemize}
\item Categorical maps preserved in the manifest for lossless cast-back.
\item Per-variable attributes kept in a sidecar record, not discarded.
\item If bidirectional: state what round-trips exactly and what is lossy.
\item Croissant: demonstrated lightly; full support is future work.
\end{itemize}
\subsection{Usage workflow}
\begin{itemize}
\item Three steps: fetch the DOI'd store, convert, point the dataloader at it.
\item Worked example with commands, runtime, and disk footprint.
\item Guidance on presets for other sampling patterns.
\end{itemize}
%% ==================================================================
%% 6  RESULTS AND BENCHMARKS
%% ==================================================================
\section{Results: the quantified trade-off}
\begin{itemize}
\item Order: throughput, sensitivity, scale-out, costs.
\item All numbers from the setup of Sect.~4.3.
\end{itemize}
\subsection{Per-sample and end-to-end throughput}
\begin{itemize}
\item $\sim$7$\times$ per-sample; $\sim$4.6$\times$ end-to-end.
\item $\sim$40$\times$ projected over the production baseline --- state assumptions.
\item Figures: latency breakdown; epoch time.
\end{itemize}
\subsection{Sensitivity to chunking strategy, store size, and worker count}
\begin{itemize}
\item Chunking: where mis-chunking erases the gain.
\item Store size: the gain grows with the store.
\item Workers: FAIR store saturates; merged store keeps scaling.
\end{itemize}
\subsection{Scale-out behaviour}
\begin{itemize}
\item Aggregate read pressure on GPFS as node count grows.
\item Does the advantage widen or narrow at scale?
\item Guidance: where each configuration becomes filesystem-limited.
\end{itemize}
\subsection{Cost of interoperability, cost of the bridge}
\begin{itemize}
\item Price of FAIR: slowdown when training directly on the interoperable store.
\item Price of the bridge: one-off conversion cost and the duplicate copy.
\item Summary table: the whole trade-off in one view.
\end{itemize}
%% ==================================================================
%% 7  CONCLUSIONS
%% ==================================================================
\conclusions  %% \conclusions[modified heading if necessary]
\begin{itemize}
\item The trade-off is real and now quantified.
\item Principle: store layout is the largest lever; metadata, not bandwidth, dominates.
\item The bridge resolves it operationally: publish FAIR, derive fast.
\item The released cube gives the community a substrate to reuse.
\item Future work: round-trip, in-place conversion, full Croissant, other filesystems.
\end{itemize}
%% The following commands are for the statements about the availability of data sets and/or software code corresponding to the manuscript.
%% It is strongly recommended to make use of these sections in case data sets and/or software code have been part of your research the article is based on.
\codeavailability{TEXT} %% use this section when having only software code available
\dataavailability{TEXT} %% use this section when having only data sets available
\codedataavailability{
%% This section discharges the data-release function of the paper. Content:
%% - Data: the 1850--2100 interoperable cube, CC BY 4.0, DOI (PANGAEA / Zenodo /
%%   BSC institutional repository — DOI must be live at submission).
%% - Code: the pipeline, the conversion (bridge) tooling, and the benchmark
%%   harness (profiler + synthetic probe), with an archived release DOI
%%   (e.g. Zenodo snapshot of the repository) as GMD requires.
%% - Exact versions used for the results in this paper.
TEXT} %% use this section when having data sets and software code available
\sampleavailability{TEXT} %% use this section when having geoscientific samples available
\videosupplement{TEXT} %% use this section when having video supplements available
\appendix
\section{}    %% Appendix A — e.g. store schemas / chunking configurations
\subsection{}     %% Appendix A1, A2, etc.
\noappendix       %% use this to mark the end of the appendix section. Otherwise the figures might be numbered incorrectly (e.g. 10 instead of 1).
%% Regarding figures and tables in appendices, the following two options are possible depending on your general handling of figures and tables in the manuscript environment:
%% Option 1: If you sorted all figures and tables into the sections of the text, please also sort the appendix figures and appendix tables into the respective appendix sections.
%% They will be correctly named automatically.
%% Option 2: If you put all figures after the reference list, please insert appendix tables and figures after the normal tables and figures.
%% To rename them correctly to A1, A2, etc., please add the following commands in front of them:
\appendixfigures  %% needs to be added in front of appendix figures
\appendixtables   %% needs to be added in front of appendix tables
%% Please add \clearpage between each table and/or figure. Further guidelines on figures and tables can be found below.
\authorcontribution{TEXT} %% this section is mandatory — settle after the lead-authorship decision (plan §6.1)
\competinginterests{TEXT} %% this section is mandatory even if you declare that no competing interests are present
\disclaimer{TEXT} %% optional section
\begin{acknowledgements}
%% CERISE belongs here only, as funding acknowledgement (grant 101082139).
TEXT
\end{acknowledgements}

\end{document}